% !TEX root = ../main.tex
%%%%%%%%%%%%%%%%%%%%%%%%%%%%%%%%%%%%%%%%%%%%%%%%%%%%%%%%%%%%%%%%%%%%%%%
% Global sensitivity analysis
%%%%%%%%%%%%%%%%%%%%%%%%%%%%%%%%%%%%%%%%%%%%%%%%%%%%%%%%%%%%%%%%%%%%%%%

\section{Global sensitivity analysis}
\label{sec:sa}

Every result so far is measured at a \emph{cell}: a chosen $\sigma$, a chosen
$c_0$, everything else at its default. We now ask whether the headline sign
survives when all sixteen defensible parameters move at once.

\subsection{The headline does not survive}

\begin{declbox}
\textbf{$P(\text{wage-led} \mid \text{viable}) = 0.095 \pm 0.007$}
(binomial standard error; 1{,}606 viable points of 3{,}328).
\textbf{Fraction viable $= 0.483$.}

\smallskip
Over the empirically defensible space, wage-led is the \emph{exception}, not
the rule --- and half of that space does not survive at all.
\end{declbox}

\paragraph{A declared design defect in the quantity of interest.}
This probability is exact about a \emph{chord} --- a two-point difference
between $\rho = 0.35$ and $\rho = 0.55$ --- and is \emph{not} the same claim as
``the derivative is negative''. \Cref{sec:frontier-fragility} had already
established that $Y(\rho)$ is U-shaped with the turning point inside the
support in 19 of 22 cells. On a U-shaped response the sign of a chord depends
on where it is taken, and can differ from an OLS slope over the full support
--- which is the estimator used in \cref{sec:wagecurve}. The sensitivity
indices correctly decompose the variance \emph{of that chord}; the viability
outcome is not a difference and is unaffected. \Cref{sec:bridge} reports the
repair.

\subsection{What actually governs the outcome}

\begin{table}[htbp]
\centering
\small
\caption{Sobol' indices. $S_1$ is the first-order index, $S_T$ the total
index; $N = 256$, bootstrap confidence intervals.}
\label{tab:sobol}
\begin{tabular}{lrrrr}
\toprule
& \multicolumn{2}{c}{Viability} & \multicolumn{2}{c}{Slope}\\
\cmidrule(lr){2-3}\cmidrule(lr){4-5}
Parameter & $S_1$ & $S_T$ & $S_1$ & $S_T$\\
\midrule
$\delta$        & \textbf{0.718} & \textbf{1.002} & \textbf{0.331} & \textbf{0.966}\\
$\pi_0$         & 0.116 & 0.290 & 0.001 & 0.562\\
$c_0$           & 0.069 & 0.086 & $-0.030$ & 0.277\\
$\lambda$       & 0.038 & 0.086 & 0.040 & 0.264\\
$\sigma$        & 0.008 & \textbf{0.008} & 0.030 & \textbf{0.024}\\
\bottomrule
\end{tabular}
\end{table}

Two facts in \cref{tab:sobol} overturn the project's own narrative.

First, $S_T \gg S_1$ throughout: the model is dominated by \emph{interactions},
as one should expect of a wage$\,\to U \to\,$capital channel that is
interactive by construction. Second, and more uncomfortably, \textbf{$\sigma$
is close to irrelevant} inside the empirical band, with a total index of
$0.024$ and a first-order index indistinguishable from zero. The depreciation
rate and the base-point capital share do the work that the project had
attributed to the elasticity of substitution.

\subsection{The depreciation cliff}
\label{sec:delta}

\begin{table}[htbp]
\centering
\small
\caption{Viability by depreciation rate.}
\label{tab:delta}
\begin{tabular}{lrrr}
\toprule
$\delta$ & Points & Fraction viable & $P(\text{wage-led} \mid \text{viable})$\\
\midrule
0.030--0.045 & 832 & \textbf{0.992} & 0.044\\
0.045--0.060 & 832 & 0.758 & 0.111\\
0.060--0.075 & 821 & 0.183 & 0.313\\
\textbf{0.075--0.090} & \textbf{843} & \textbf{0.000} & ---\\
\bottomrule
\end{tabular}
\end{table}

$\delta$ alone governs survival: its total index on viability is $1.00$, and no
other parameter approaches it. \Cref{tab:delta} is not a robustness note but a
structural limitation, and it becomes one only when combined with the anchoring
exercise. The depreciation rate \emph{implied by the data for this model's
perimeter} --- private non-residential fixed capital on both sides of the
ratio --- is $\approx 0.090$. In that band, \textbf{zero of 843 sampled points
survive}.

\begin{declbox}[frametitle={Why this cannot be recalibrated away}]
The tempting reading is that the analysis \emph{confirms} $\delta = 0.05$. It
does the opposite, and the reason is the steady-state accounting identity
$I/K = \delta + g$. The data satisfy it with $g \approx 0.022$; this model has
$g = 0$, because technical change is out of scope. With $g = 0$, steady-state
investment must cover \emph{all} depreciation and nothing else. Raising
$\delta$ toward the data without adding growth asks the capital block to
finance a growing economy's depreciation out of a stationary economy's
investment flow: capital erodes every period and the economy dies.

$\delta = 0.05$ is therefore \emph{not} where the data put it. It is where
$g = 0$ forces it to sit for the model to exist at all. The repair is a
mechanism --- trend productivity growth --- not a calibration.
\end{declbox}

This is the same structural signature as the model's other declared failure of
fit. Measured $I/Y$ ($0.158$--$0.182$ at $\rho = 0.40$) sits above the BEA
anchor of $0.138$--$0.141$, and measured $K/Y$ ($3.17$--$3.64$) sits far above
the data's $1.23$. The second follows mechanically from the first: with
$g = 0$, $I = \delta K$ and hence $K/Y = (I/Y)/\delta$; measured $I/K$ is
$0.0500$, exactly $\delta$. \emph{The model cannot match business $I/Y$ and
$K/Y$ simultaneously without growth.} Seen from the level side, $K/Y$ does not
fit; seen from the viability side, the empirical $\delta$ kills the economy.
One cause, two symptoms.

\subsection{The wide-\texorpdfstring{$\sigma$}{sigma} check, and an inverted
direction}
\label{sec:inversion}

Widening the $\sigma$ range to $0.30$--$1.00$ leaves viability identical
($0.483$) but \emph{doubles} the wage-led probability to $0.201$, and the
change is sharply concentrated (\cref{tab:widesigma}).

\begin{table}[htbp]
\centering
\small
\caption{Wide-$\sigma$ check ($N = 128$, declared secondary).}
\label{tab:widesigma}
\begin{tabular}{lrrr}
\toprule
$\sigma$ & Fraction viable & $P(\text{wage-led} \mid \text{viable})$
& Mean slope\\
\midrule
0.301--0.475 & 0.519 & 0.042 & $+53.7$\\
0.475--0.649 & 0.464 & 0.057 & $+46.7$\\
\textbf{0.649--0.823} & 0.469 & \textbf{0.338} & $+27.2$\\
\textbf{0.823--0.998} & 0.481 & \textbf{0.380} & $+21.5$\\
\bottomrule
\end{tabular}
\end{table}

The pattern is monotone in $\sigma$ and in the \emph{same} direction as
\cref{sec:frontier}: \cref{sec:frontier} puts the empirical range below
$\sigstar$, where the model is not wage-led, and here too wage-led is rare
below $0.65$ and commoner above it. The two are consistent.

\begin{declbox}[frametitle={A reading error, corrected in \cref{sec:bridge}}]
An earlier draft of this section read the table above as showing a direction
\emph{inverted} relative to the conditional result. That reading was wrong, and
it was wrong twice over. It is recorded here rather than deleted, because the
experiment of \cref{sec:bridge} was designed to explain it and the design
survives its own premise.

First, the direction was never inverted: ``the empirical range sits below
$\sigstar$ and is therefore not wage-led'' and ``below $0.65$ wage-led is rare''
are the same statement, not opposite ones. Second, there is no threshold at
$0.65$ --- or anywhere: $P(\text{wage-led} \mid \text{viable})$ never reaches
$0.5$ at any $\sigma$ tested, topping out at $0.380$. \Cref{sec:bridge} measures
the marginalised frontier directly and puts it at $\sigstar \approx 0.96$.
\end{declbox}

What \emph{is} genuinely at issue is the estimator. The chord used here and the
OLS slope used in \cref{sec:frontier} are different functionals of a curved
$Y(\rho)$, and we deliberately do not resolve which to prefer by argument. Two
causes are candidates:

\begin{enumerate}[label=(\alph*)]
  \item \textbf{Chord versus derivative.} \Cref{sec:wagecurve} fits an OLS
  slope over the whole support; \cref{sec:sa} measures a chord over a
  sub-interval. On a U-shaped $Y(\rho)$ with the turning point inside the
  support, the two can differ in sign with neither being wrong.
  \item \textbf{Conditional versus marginal.} $\sigstar$ is a
  \emph{conditional} statement, measured with every other parameter at its
  default. The sensitivity analysis measures a \emph{marginal} effect,
  averaging over draws of the other fifteen. In a model where $S_T \gg S_1$
  there is no reason these should agree --- and this hypothesis already has
  independent support in \cref{tab:sobol}, where marginally $\sigma$ explains
  about 2\% of the variance. A robust conditional $\sigstar$ and a nearly inert
  marginal $\sigma$ are entirely compatible.
\end{enumerate}

\Cref{fig:sa} shows the screening step and the variance decomposition that
underlie \cref{tab:sobol,tab:delta}.

\begin{figure}[htbp]
\centering
\begin{subfigure}{0.48\textwidth}
  \includegraphics[width=\linewidth]{ces_b13_morris_mu_sigma.png}
  \caption{Morris screening.}
\end{subfigure}\hfill
\begin{subfigure}{0.48\textwidth}
  \includegraphics[width=\linewidth]{ces_b13_sobol_indices.png}
  \caption{Sobol' indices with bootstrap CIs.}
\end{subfigure}
\caption{Screening and decomposition over the sixteen-parameter space.}
\label{fig:sa}
\end{figure}

\subsection{Two by-products}

\paragraph{Kalecki, confirmed in levels.} Within this model the relation
$\Pi - I = C - W$ is a tautology and cannot settle causation. Sobol' sampling,
however, varies the capitalists' propensity to consume \emph{independently of
everything else}, which makes the comparison interventional. Moving from the
lowest to the highest quartile of capitalist MPC raises capitalist consumption
by $10.83$ and profit by $11.56$ --- near one-for-one, with a correlation of
$+0.83$ on levels. Capitalists do earn what they spend. The profit
\emph{share} falls slightly, because output rises faster; Kalecki's claim is
about levels, and it is the levels that hold.

\paragraph{The accelerator produces the sign.} Wage-led outcomes are
overwhelmingly concentrated at high accelerator reactivity. In the lowest
quartile of $\beta$ ($[0.002, 0.251]$) only $0.6\%$ of viable points are
wage-led (2 of 353); in the highest ($[0.750, 0.999]$) the figure is $22.1\%$,
and viability itself is higher ($0.523$ against $0.423$). The wage-led sign is
therefore substantially \emph{produced by} investment reactivity rather than
being an intrinsic property of the technology, and $\beta$ --- a parameter
with no canonical empirical value (\cref{tab:params}) --- governs both the
sign and survival. This reverses the expectation that the instability would
prove independent of investment reactivity, and it is an uncomfortable result:
the headline depends on the one parameter in the model with the weakest
empirical claim.
